% Generated draft: requires mathematical and source review.
\documentclass[11pt]{article}
\usepackage[margin=25mm]{geometry}
\usepackage{fontspec}
\setmainfont{Latin Modern Roman}
\usepackage{amsmath,amssymb,mathtools,unicode-math}
\setmathfont{Latin Modern Math}
\usepackage{enumitem,needspace,xurl,hyperref}
\setlength{\parindent}{0pt}
\setlength{\parskip}{5pt}
\setlength{\emergencystretch}{3em}
\newcommand{\N}{\mathbb{N}}
\newcommand{\Z}{\mathbb{Z}}
\newcommand{\Q}{\mathbb{Q}}
\newcommand{\R}{\mathbb{R}}
\newcommand{\C}{\mathbb{C}}
\newcommand{\F}{\mathbb{F}}
\newcommand{\A}{\mathbb{A}}
\newcommand{\PP}{\mathbb{P}}
\newcommand{\E}{\mathbb{E}}
\newcommand{\eps}{\varepsilon}
\newcommand{\dd}{\,\mathrm{d}}
\begin{document}
\section*{UnsolvedMath ID: 30003801}
\textbf{Source problem code:} catalogue entry
\subsection{Definitions and mathematical statement}
Let \(F\) denote the common base field, and let \(q:V\to F\) and \(q':V'\to F\) be quadratic forms with \(\dim_F V=\dim_F V'\). The supplied record does not specify the characteristic of \(F\), the nondegeneracy hypotheses, or the coefficient ring for Chow motives. Consequently, the following is a conditional formulation whose admissible setting still requires clarification.

Write
\[
Q(q)=\{[v]\in\mathbb P(V):q(v)=0\}
\]
for the projective quadric of \(q\), and similarly for \(q'\). For smooth projective quadrics and a fixed coefficient ring \(\Lambda\), write \(M(Q(q);\Lambda)\) for the Chow motive. The ring \(\Lambda\) is a placeholder for a choice absent from the source, not an assertion that the question is intended for every coefficient ring.

For an extension field \(E/F\), put \(V_E=V\otimes_F E\) and denote the extended form by \(q_E\). In the usual nonsingular setting, its Witt index is
\[
i_W(q_E)=\max\{\dim_E U:U\subseteq V_E\text{ is an }E\text{-linear subspace and }q_E|_U=0\}.
\]
The forms \(q\) and \(q'\) are similar over \(F\) if there exist \(\lambda\in F^\times\) and an \(F\)-linear isomorphism \(f:V\to V'\) such that \(q'(f(v))=\lambda q(v)\) for every \(v\in V\).

The first question in the supplied record is whether, under the premise
\[
M(Q(q);\Lambda)\simeq M(Q(q');\Lambda),
\]
one has
\[
\bigl[\forall\text{ extension fields }E/F,\quad i_W(q_E)=i_W(q'_E)\bigr]
\quad\Longleftrightarrow\quad
\bigl[q\text{ and }q'\text{ are similar over }F\bigr].
\]
The quantifier over extension fields is unrestricted in the supplied statement. This question is recorded without asserting either implication or supplying the missing hypotheses.

For an integer \(k\) with \(1\leq k\leq\dim_F V\), let \(\operatorname{OGr}_k(q)\) denote the variety of totally \(q\)-isotropic \(k\)-dimensional subspaces of \(V\); its \(E\)-points correspond to such subspaces of \(V_E\). The second question asks whether the motivic criterion can instead use motives of these isotropic Grassmannians. The record does not specify \(k\), whether one uses one Grassmannian or a family, whether connected components are distinguished, or exactly which equivalence is to be preserved.

Finally, for a semisimple algebraic group \(G\) over \(F\), a projective homogeneous \(G\)-variety means a projective \(F\)-variety with a \(G\)-action that becomes transitive after extension to an algebraic closure. The source asks for an extension of the motivic criterion to all semisimple groups using motives of such varieties. It does not specify the comparison relation between groups, the required homogeneous varieties, their types, the motive coefficients, or the invariant that would replace similarity or Witt indices. Thus this final request is a research direction rather than a uniquely determined quantified assertion.
\subsection{Short English statement}
The supplied question compares isomorphic quadric motives, equality of Witt indices over every field extension, and similarity of quadratic forms. It also asks about replacing quadrics by isotropic Grassmannians and extending the comparison to semisimple groups. Essential hypotheses and the precise extension targets remain unspecified. The supplied status is \emph{partially solved}; this draft does not independently verify that assessment.
\subsection{Sources}
\begin{itemize}
\item Local UnsolvedMath registry: permanent ID \(30003801\), title \emph{Motivic Equivalence and Similarity of Quadratic Forms}, collection \emph{Oberwolfach Reports Open Problems}, exact source problem code \texttt{catalogue entry}. The supplied record lists proposed year 2018 and names no proposer.
\item Supplied report citation: \emph{Quadratic Forms and Related Structures over Fields (2019)}, Oberwolfach Reports, pp.\ 1235--1291, DOI \texttt{10.4171/owr/2018/21}. \url{https://doi.org/10.4171/owr/2018/21}. The supplied background describes these reports as records of mathematical workshops and their open-problem discussions. No problem-specific page or passage is supplied.
\item Supplied literature reference, labelled ``De Clercq on motivic equivalence'': \url{https://doi.org/10.1017/S0010437X17000644}.
\item Supplied literature reference, labelled ``Critical varieties of semisimple groups'': \url{https://arxiv.org/abs/1802.03857}.
\item Supplied literature assessment dated \texttt{2026-08-21}: status \emph{partially solved}. The assessment says that the quadratic form cases and the negative Grassmannian extension are established, while motivic equivalence and critical varieties are not classified for all semisimple groups. This is supplied dated curation derived from the OWR material, not an independently verified literature assessment; time-sensitive status and cited claims require independent verification.
\item The supplied registry labels the statement \texttt{corrected\_verified} and reports a manual semantic, typographic, and LaTeX review that deemed it self-contained and faithful to its target. These are supplied editorial claims, not mathematical verification by this draft.
\end{itemize}
\end{document}
